\section{Thiết kế hệ thống}

\subsection{Thiết kế kiến trúc hệ thống}

\subsubsection{Tổng quan kiến trúc}

Hệ thống Micro-learning được thiết kế theo kiến trúc \textbf{Hexagonal
(Port/Adapter)} do Alistair Cockburn đề xuất, tổ chức mã nguồn thành ba tầng
chính: \textbf{Domain} (nghiệp vụ cốt lõi + Port interface),
\textbf{Application} (use case điều phối), và \textbf{Adapter} (hiện thực
hạ tầng kỹ thuật). Hai tầng giao tiếp bên ngoài là \textbf{Delivery}
(gRPC + HTTP API + BullMQ Workers) và \textbf{Infrastructure} (PostgreSQL,
Qdrant, Neo4j, MinIO, Redis, Gemini LLM).

Toàn bộ backend được đóng gói trong một Python monolith có cấu trúc chặt
chẽ thay vì nhiều microservice riêng biệt — đây là lựa chọn có chủ đích
cho giai đoạn đồ án: giảm độ phức tạp triển khai, tận dụng ưu điểm của
Python import system để thực thi dependency injection qua Container, và
dễ dàng kiểm thử toàn bộ pipeline mà không cần mock network.

\begin{figure}[H]
    \centering
    \includegraphics[width=\textwidth]{figures/overview.png}
    \caption{Kiến trúc tổng quan: AI Backend (chunking\_v2) +
             Open edX (Redwood) + Next.js Frontend qua mạng Docker
             \texttt{lms\_shared}.}
    \label{fig:arch-overview}
\end{figure}

Hệ thống gồm bốn nhóm thành phần chính:

\begin{itemize}
    \item \textbf{AI Backend (chunking\_v2):} Python monolith tổ chức
          theo Hexagonal Architecture, expose 12 gRPC RPCs qua proto3
          và HTTP API qua FastAPI. Ba BullMQ worker xử lý pipeline bất
          đồng bộ (document, enrichment, outbox).
    \item \textbf{Open edX LMS (Redwood):} Tutor 18.2.2, tích hợp
          LTI 1.3 qua \texttt{lti\_consumer-xblock} có sẵn. Custom
          \texttt{AILessonBlock} đóng vai trò ``thin glue'' — chỉ
          render iframe trỏ đến Next.js frontend.
    \item \textbf{Next.js Frontend:} Next.js 15 App Router +
          TypeScript + Tailwind, xác thực LTI 1.3 qua \texttt{ltijs}
          Provider, gọi gRPC server-side qua proxy route, quản lý
          session bằng Redis (TTL 30 phút).
    \item \textbf{Infrastructure:} PostgreSQL (14 bảng metadata +
          curriculum), Qdrant (vector similarity), Neo4j (Knowledge
          Graph Chuẩn đầu ra), MinIO (object storage), Redis (cache +
          session + job queue).
\end{itemize}

\subsubsection{Kiến trúc phân tầng Hexagonal}

Hình~\ref{fig:arch-layers} minh họa kiến trúc ba tầng kèm tầng
Delivery và Adapter ở hai phía. Mũi tên phụ thuộc luôn hướng vào
tầng Domain --- các tầng ngoài phụ thuộc vào Port (interface), Domain
không biết gì về Adapter cụ thể.

\begin{figure}[H]
    \centering
    \includegraphics[width=\textwidth]{figures/arch_hexagonal.png}
    \caption{Kiến trúc Hexagonal của chunking\_v2: Domain core ở
             trung tâm, Application điều phối, Adapter hiện thực hạ tầng.}
    \label{fig:arch-layers}
\end{figure}

\begin{table}[H]
    \centering
    \caption{Mô tả ba tầng kiến trúc Hexagonal của chunking\_v2.}
    \label{tab:arch-layers}
    \begin{tabular}{p{3cm}p{4.5cm}p{7cm}}
        \toprule
        \textbf{Tầng} & \textbf{Thành phần} & \textbf{Trách nhiệm} \\
        \midrule
        \textbf{Domain}
            & Entities, Ports (ABC), Exceptions, Result[T,E]
            & Logic nghiệp vụ thuần túy. Không import bất kỳ
              framework hay adapter nào. Mọi dependency được
              định nghĩa qua Port interface. \\
        \midrule
        \textbf{Application}
            & 12 Use Cases (ProcessDocument, IngestCurriculum,
              GenerateCurriculumQuiz, \dots)
            & Điều phối domain entities qua Port interface.
              Mỗi use case nhận input DTO, gọi domain + port,
              trả Result[T,E]. \\
        \midrule
        \textbf{Adapter}
            & PostgresMetadataStore, QdrantVectorStore,
              Neo4jGraphStore, MinIOAdapter, GeminiAdapter,
              BGEEmbedder, DcmhExtractor, \dots
            & Hiện thực các Port interface với công nghệ cụ thể.
              Có thể swap Adapter (vd: Gemini $\to$ OpenAI) mà
              không sửa tầng Application hay Domain. \\
        \bottomrule
    \end{tabular}
\end{table}

\textbf{Delivery Layer} (primary adapter): gRPC Server
(\texttt{chunking\_pb2\_grpc}), HTTP API (FastAPI), BullMQ Workers
(document, enrichment, outbox), và Next.js Server (LTI handler).
Tất cả đều gọi Use Case thông qua dependency injection container.

\subsubsection{Bài toán thiết kế và giải pháp}

Nhóm xác định năm bài toán kỹ thuật chính:

\begin{table}[H]
    \centering
    \caption{Năm bài toán thiết kế và giải pháp.}
    \label{tab:design-problems}
    \begin{tabular}{p{0.5cm}p{4cm}p{4cm}p{5.5cm}}
        \toprule
        \textbf{\#} & \textbf{Bài toán} & \textbf{Nguyên nhân} & \textbf{Giải pháp} \\
        \midrule
        1 & Xử lý bất đồng bộ
          & AI pipeline kéo dài vài phút; HTTP/gRPC sẽ timeout.
          & BullMQ job queue + 3 workers + state machine. \\
        \midrule
        2 & Thay thế LLM provider
          & Gemini có thể thay đổi API; cần linh hoạt.
          & Port/Adapter (ILLMClient $\to$ GeminiAdapter). \\
        \midrule
        3 & Kiểm thử không gọi API thật
          & Gọi LLM API thật tốn tiền và phụ thuộc mạng.
          & Mock Adapter implement cùng Port interface. \\
        \midrule
        4 & Chatbot trả lời sai kiến thức
          & LLM có thể sinh thông tin không có trong bài học.
          & RAG --- chỉ trả lời dựa trên context từ Qdrant. \\
        \midrule
        5 & Tích hợp LMS qua LTI 1.3
          & Open edX có sẵn LTI consumer; cần thin client ở Next.js.
          & Thin XBlock glue + ltijs Provider + OIDC flow. \\
        \bottomrule
    \end{tabular}
\end{table}

\paragraph{Bài toán 1 --- Job Queue + State Machine.}

Khi giảng viên upload tài liệu, gRPC handler \texttt{EnqueueDocument}
đẩy job vào Redis-backed BullMQ queue và trả \texttt{document\_id}
ngay lập tức. Ba worker riêng biệt tiêu thụ queue:

\begin{itemize}
    \item \texttt{document-worker}: parse $\to$ chunk $\to$ embed
          $\to$ upsert Qdrant + Postgres.
    \item \texttt{enrichment-worker}: gọi Gemini sinh card + quiz
          từ chunk đã embed.
    \item \texttt{outbox-worker}: đồng bộ Postgres $\to$ Neo4j
          Knowledge Graph qua event sourcing pattern.
\end{itemize}

Pipeline theo dõi trạng thái qua state machine với các trạng thái:
\texttt{QUEUED} $\to$ \texttt{PARSING} $\to$ \texttt{CHUNKING}
$\to$ \texttt{EMBEDDING} $\to$ \texttt{ENRICHING} $\to$
\texttt{DONE} (hoặc \texttt{ERROR}). Cơ chế retry 3 lần với
exponential backoff ($1s \to 2s \to 4s$).

\paragraph{Bài toán 2 \& 3 --- Hexagonal Architecture.}

Mỗi dependency bên ngoài được trừu tượng hóa thành một \textbf{Port}
(abstract base class trong \texttt{src/domain/ports/}). Khi cần
thay thế provider, chỉ cần tạo Adapter mới implement cùng Port mà
không sửa use case hay domain entity nào.

\begin{table}[H]
    \centering
    \caption{Danh sách Port và Adapter trong hệ thống.}
    \label{tab:ports-adapters}
    \begin{tabular}{p{3.5cm}p{5cm}p{6cm}}
        \toprule
        \textbf{Port (Interface)} & \textbf{Production Adapter} & \textbf{Test Adapter} \\
        \midrule
        \texttt{IParser}
            & PdfParser, DocxParser, PptxParser
            & MockParser \\
        \texttt{IEmbeddingModel}
            & BGEEmbedder (BGE-M3, 1024-dim, local)
            & MockEmbeddingAdapter \\
        \texttt{IVectorStore}
            & QdrantVectorStore
            & InMemoryVectorStore \\
        \texttt{IMetadataStore}
            & PostgresMetadataStore (14 tables)
            & MockMetadataStore \\
        \texttt{IGraphStore}
            & Neo4jGraphStore (LO Knowledge Graph)
            & MockGraphStore \\
        \texttt{IObjectStorage}
            & MinIOAdapter
            & LocalFSAdapter \\
        \texttt{ILLMClient}
            & GeminiAdapter (gemini-3-flash)
            & MockLLMClient \\
        \texttt{ICurriculumExtractor}
            & DcmhExtractor (regex CDIO + Bloom)
            & --- (structural, không cần mock) \\
        \texttt{IChunkLOMapper}
            & HeuristicLoMapper (heading match)
            & --- \\
        \bottomrule
    \end{tabular}
\end{table}

\paragraph{Bài toán 4 --- Curriculum-Aware Design.}

Hệ thống thiết kế lớp curriculum-aware xuyên suốt pipeline:

\begin{enumerate}
    \item \textbf{Extract:} \texttt{DcmhExtractor} parse DCMH PDF
          bằng regex CDIO pattern (\texttt{L.O.X.Y}) và ánh xạ
          động từ Bloom sang 6 mức (Remember $\to$ Create).
    \item \textbf{Persist:} Curriculum được lưu vào PostgreSQL
          (6 bảng: courses, chapters, learning\_outcomes,
          assessments, lo\_assessments, chunk\_lo\_mappings)
          và Neo4j Knowledge Graph (node LO, Chapter, Assessment;
          edge HAS\_CHAPTER, HAS\_LO, ASSESSES, SUPPORTS).
    \item \textbf{Map:} \texttt{HeuristicLoMapper} tự động liên kết
          chunk tài liệu với LO dựa trên heading alignment.
    \item \textbf{Generate:} \texttt{GenerateCurriculumQuiz} gọi
          Gemini sinh câu hỏi MCQ + short answer \textbf{aligned}
          với từng LO, kèm Bloom level và rationale.
\end{enumerate}

\paragraph{Bài toán 5 --- Tích hợp LMS qua LTI 1.3.}

Quyết định sử dụng LTI 1.3 (OAuth2 + JWT) thay vì LTI 1.1
(OAuth1 + XML) vì tính bảo mật cao hơn và hỗ trợ AGS
(Assignment and Grade Services) cho grade passback tự động.

XBlock \texttt{AILessonBlock} đóng vai trò ``thin glue'' ---
không chứa logic nghiệp vụ, chỉ render iframe với OIDC login URL
trỏ đến Next.js. Toàn bộ logic xác thực và hiển thị nằm ở Next.js.

\begin{figure}[H]
    \centering
    \includegraphics[width=\textwidth]{figures/lti_auth_flow.png}
    \caption{Luồng xác thực LTI 1.3 OIDC: từ click XBlock đến
             /learn/cards với session Redis 30 phút.}
    \label{fig:lti-auth-flow}
\end{figure}

Luồng xác thực (hình~\ref{fig:lti-auth-flow}):
\begin{enumerate}
    \item Sinh viên đăng nhập LMS, vào unit chứa AILessonBlock.
    \item XBlock render iframe với OIDC login URL trỏ Next.js.
    \item Next.js \texttt{/lti/login} khởi tạo OIDC flow:
          redirect browser đến edX \texttt{/oauth2/authorize}.
    \item edX xác thực sinh viên, sinh \texttt{id\_token} JWT,
          auto-submit form POST đến Next.js \texttt{/lti/launch}.
    \item \texttt{ltijs} Provider validate JWT qua edX JWKS,
          extract claims: \texttt{sub}, \texttt{context.id},
          \texttt{roles}, \texttt{custom}.
    \item Postgres: UPSERT \texttt{lms\_user\_mappings} tạo
          \texttt{internal\_user\_id}.
    \item Redis: SET session với TTL 30 phút.
    \item Set HttpOnly cookie \texttt{sid} và 302 redirect:
          giảng viên $\to$ \texttt{/manage/review},
          sinh viên $\to$ \texttt{/learn/cards}.
\end{enumerate}

\subsubsection{Biện luận kiến trúc: Polyglot Monolith cho nhóm 4 thành viên}

\textbf{Tổ chức nhóm theo kiến trúc.} Dự án được thực hiện bởi nhóm 4
thành viên trong 30 tuần. Kiến trúc Hexagonal Monolith được chọn thay vì
Microservices vì ba lý do chính:

\begin{enumerate}
    \item \textbf{Không cần network isolation:} Tất cả worker (document,
          enrichment, outbox --- và GĐ2 thêm video, chat) chia sẻ chung
          Domain entities và Adapter (PostgreSQL, Qdrant, Neo4j).
          Microservices sẽ buộc mỗi service maintain schema riêng + API
          contract riêng, tăng chi phí đồng bộ giữa 4 người.
    \item \textbf{Phân công không đụng hàng:} Mỗi thành viên làm việc
          trong một thư mục độc lập trong cùng monolith:
          \texttt{src/adapters/curriculum/} (Thành viên A),
          \texttt{src/adapters/video/} (Thành viên B),
          \texttt{src/adapters/tracer/} (Thành viên C),
          \texttt{web/} Next.js (Thành viên D).
          Domain Ports (interface) được thiết kế trước --- mỗi người
          implement Adapter theo Port, merge không conflict.
    \item \textbf{Dễ kiểm thử tích hợp:} Toàn bộ pipeline test được
          trong 1 process với Mock Adapter, không cần setup multi-service
          test environment.
\end{enumerate}

\textbf{Chiến lược Polyglot.} Hệ thống sử dụng ba ngôn ngữ, mỗi ngôn
ngữ đảm nhận domain mà nó mạnh nhất:

\begin{table}[H]
    \centering
    \caption{Phân bổ ngôn ngữ theo domain.}
    \label{tab:polyglot}
    \begin{tabular}{p{2.5cm}p{5cm}p{7cm}}
        \toprule
        \textbf{Ngôn ngữ} & \textbf{Thành phần} & \textbf{Lý do chọn} \\\\
        \midrule
        \textbf{Python 3.10+}
            & chunking\_v2 backend (gRPC server, workers, adapters)
            & Hệ sinh thái ML/AI mạnh nhất: \texttt{sentence-transformers}
              (BGE-M3), \texttt{faster-whisper}, \texttt{pdfplumber},
              \texttt{neo4j}, \texttt{psycopg}, \texttt{grpcio}.
              Hexagonal Architecture với ABC + dataclass rất tự nhiên
              trong Python. \\\\
        \textbf{TypeScript (Node.js)}
            & Next.js 15 frontend (LTI auth, UI, gRPC proxy)
            & \texttt{ltijs} là thư viện LTI 1.3 mạnh nhất hiện có,
              chỉ có bản Node.js. Next.js App Router cho SSR + API route
              trong cùng process. \\\\
        \textbf{Go (Golang)}
            & Video Processing Worker (GĐ2)
            & Goroutines xử lý concurrent video chunk encoding/decoding
              hiệu quả hơn Python asyncio. Hệ sinh thái FFmpeg binding
              (\texttt{go-ffmpeg}) mạnh. Giao tiếp với backend qua
              gRPC (\texttt{protobuf} native trong Go). \\
        \bottomrule
    \end{tabular}
\end{table}

Video Worker bằng Go giao tiếp với chunking\_v2 Python backend
qua gRPC --- đây cũng là ví dụ thực tế cho thấy Port/Adapter pattern
hoạt động xuyên ngôn ngữ: Go service gọi \texttt{EnqueueVideo} RPC
để nhận job, và gọi \texttt{UpdateDocumentStatus} để báo tiến độ
về Python backend.

\textbf{Tóm lại:} Kiến trúc Hexagonal Monolith + Polyglot Workers
là lựa chọn tối ưu cho đồ án 4 người --- tận dụng ưu điểm từng ngôn
ngữ, giữ chi phí tích hợp thấp, và vẫn đảm bảo khả năng mở rộng
thành microservices sau này nếu cần.

\subsection{Thiết kế Domain Model}

\subsubsection{Entities cốt lõi}

Domain model của hệ thống được thiết kế quanh các entity chính sau:

\begin{table}[H]
    \centering
    \caption{Các entity chính trong Domain Model.}
    \label{tab:domain-entities}
    \begin{tabular}{p{3.5cm}p{8cm}p{3cm}}
        \toprule
        \textbf{Entity} & \textbf{Mô tả} & \textbf{File} \\
        \midrule
        \texttt{Document}
            & Tài liệu gốc (PDF, DOCX, PPTX). Chứa metadata:
              document\_id, name, doc\_type, mime\_type,
              course\_id, owner\_id, language.
            & \texttt{document.py} \\
        \texttt{ParsedDocument}
            & Kết quả parse: danh sách Section (heading,
              content text, images, page\_number).
            & \texttt{document.py} \\
        \texttt{Chunk}
            & Đoạn văn bản đã chia theo heading, kèm
              heading\_path, page\_number, vector embedding.
            & \texttt{chunk.py} \\
        \texttt{LessonCard}
            & Thẻ bài học: title, bullets (3--5 ý), key\_insight,
              source\_chunk\_ids, card\_index.
            & \texttt{card.py} \\
        \texttt{QuizItem}
            & Câu hỏi: question, choices[], correct\_index,
              explanation, difficulty, lo\_id, bloom\_level.
            & \texttt{quiz.py} \\
        \texttt{Course}
            & Khóa học: code, title\_vi, credits,
              extraction\_confidence từ DCMH.
            & \texttt{curriculum.py} \\
        \texttt{Chapter}
            & Chương trong DCMH: code, title, order\_index.
            & \texttt{curriculum.py} \\
        \texttt{LearningOutcome}
            & Chuẩn đầu ra: code (L.O.X.Y), statement\_vi,
              bloom\_level, cdio\_level.
            & \texttt{curriculum.py} \\
        \texttt{Assessment}
            & Bài đánh giá: code, name\_vi, category,
              weight (0--1).
            & \texttt{curriculum.py} \\
        \bottomrule
    \end{tabular}
\end{table}

\subsubsection{Port Interface (Abstract Base Classes)}

Tất cả dependency bên ngoài được định nghĩa qua ABC trong
\texttt{src/domain/ports/}:

\begin{itemize}
    \item \texttt{IParser} --- parse file $\to$ ParsedDocument.
    \item \texttt{IChunker} --- chia ParsedDocument $\to$ List[Chunk].
    \item \texttt{IEmbeddingModel} --- embed text $\to$ vector[1024].
    \item \texttt{IVectorStore} --- upsert/search/delete vectors.
    \item \texttt{IMetadataStore} --- CRUD document/chunk/card/quiz.
    \item \texttt{IGraphStore} --- Neo4j curriculum graph operations.
    \item \texttt{IObjectStorage} --- upload/download file (MinIO).
    \item \texttt{ILLMClient} --- gọi LLM sinh text (Gemini).
    \item \texttt{ICurriculumExtractor} --- extract LO từ DCMH PDF.
    \item \texttt{IChunkLOMapper} --- map chunk $\to$ LO.
\end{itemize}

\subsubsection{Result Monad}

Hệ thống dùng \texttt{Result[T, E]} monad để xử lý lỗi không
dùng exception trong tầng domain. Mọi Port method trả về
\texttt{Result[Ok[T], Err[Exception]]}:

\begin{lstlisting}[language=Python, caption={Result monad trong domain}]
from src.shared.result import Ok, Err, Result

def process_document(...) -> Result[Document, Exception]:
    if not file_data:
        return Err(ValueError("empty file"))
    ...
    return Ok(document)
\end{lstlisting}

\subsection{Thiết kế Use Case Layer}

Mỗi gRPC RPC trong \texttt{chunking.proto} được ánh xạ 1--1
với một Use Case class trong \texttt{src/application/}:

\begin{table}[H]
    \centering
    \caption{12 RPCs và Use Case tương ứng.}
    \label{tab:rpc-usecases}
    \begin{tabular}{p{5cm}p{5.5cm}p{4cm}}
        \toprule
        \textbf{gRPC RPC} & \textbf{Use Case} & \textbf{Loại} \\
        \midrule
        \texttt{HealthCheck}
            & HealthCheckUseCase
            & Đồng bộ \\
        \texttt{ProcessDocument}
            & ProcessDocumentUseCase
            & Đồng bộ (nặng) \\
        \texttt{EnqueueDocument}
            & EnqueueDocumentUseCase
            & Bất đồng bộ \\
        \texttt{GetDocumentStatus}
            & GetDocumentStatusUseCase
            & Polling \\
        \texttt{GetCards}
            & GetCardsUseCase
            & Đọc \\
        \texttt{GetQuiz}
            & GetQuizUseCase
            & Đọc \\
        \texttt{Search}
            & SearchChunksUseCase
            & Vector search \\
        \texttt{DeleteDocument}
            & DeleteDocumentUseCase
            & Xóa \\
        \texttt{IngestCurriculum}
            & IngestCurriculumUseCase
            & Extract + Persist \\
        \texttt{GetCurriculum}
            & GetCurriculumUseCase
            & Đọc \\
        \texttt{SearchByLearningOutcome}
            & SearchByLearningOutcomeUseCase
            & LO-scoped search \\
        \texttt{GenerateCurriculumQuiz}
            & GenerateCurriculumQuizUseCase
            & LO-aligned quiz \\
        \bottomrule
    \end{tabular}
\end{table}

Sơ đồ tuần tự (sequence diagram) cho luồng \texttt{IngestCurriculum}
được minh họa trong hình~\ref{fig:seq-ingest-curriculum}.

\begin{figure}[H]
    \centering
    \includegraphics[width=\textwidth]{figures/seq_ingest_curriculum.png}
    \caption{Sơ đồ tuần tự: IngestCurriculum --- parse DCMH PDF
             $\to$ extract LO $\to$ persist Postgres + Neo4j.}
    \label{fig:seq-ingest-curriculum}
\end{figure}

\subsection{Thiết kế Adapter Layer}

\subsubsection{PostgreSQL Metadata Store}

Adapter \texttt{PostgresMetadataStore} implement
\texttt{IMetadataStore} với 14 bảng:

\begin{itemize}
    \item \textbf{Core (6):} documents\_metadata, chunks\_metadata,
          chunk\_contents, lesson\_cards, quiz\_items, outbox\_events.
    \item \textbf{Curriculum (6):} courses, chapters, learning\_outcomes,
          assessments, lo\_assessments, chunk\_lo\_mappings.
    \item \textbf{LMS Integration (2):} lms\_user\_mappings,
          lms\_course\_ref.
    \item \textbf{Concepts (2):} concepts, chunk\_concepts.
\end{itemize}

Schema được tự động tạo qua \texttt{\_ensure\_schema()} với cơ chế
double-checked locking đảm bảo thread-safe initialization.

\subsubsection{Qdrant Vector Store}

Collection \texttt{documents}: vector 1024 chiều (BGE-M3), metric
cosine. Payload gồm document\_id, chunk\_id, heading\_path,
language, course\_id. Hỗ trợ filter theo course\_id, doc\_type,
language để giới hạn phạm vi tìm kiếm.

\subsubsection{Neo4j Knowledge Graph}

Schema graph gồm các node label \texttt{Course}, \texttt{Chapter},
\texttt{LearningOutcome}, \texttt{Assessment}, \texttt{Chunk},
\texttt{Concept} và các relationship \texttt{HAS\_CHAPTER},
\texttt{HAS\_LO}, \texttt{ASSESSES}, \texttt{SUPPORTS},
\texttt{RELATED\_TO}. Dữ liệu được đồng bộ từ PostgreSQL qua
Outbox Worker theo pattern event-driven.

Bốn method chính đã được hiện thực:
\texttt{upsert\_curriculum\_graph},
\texttt{upsert\_chunk\_lo\_mappings},
\texttt{find\_chunks\_for\_lo},
\texttt{find\_chunks\_for\_chapter}.

\subsubsection{Gemini LLM Adapter}

Adapter \texttt{GeminiAdapter} implement \texttt{ILLMClient} với
model \texttt{gemini-3-flash-preview}. Prompt template cho từng
loại tác vụ:

\begin{itemize}
    \item \textbf{Card generation:} ``Sinh 3--5 ý chính dạng bullet
          points từ đoạn văn sau. Mỗi ý không quá 30 từ.''
    \item \textbf{Quiz generation:} ``Sinh 5 câu MCQ + 2 câu trả
          lời ngắn. Mỗi câu có 4 lựa chọn, 1 đáp án đúng, kèm
          giải thích.''
    \item \textbf{Curriculum quiz:} ``Sinh câu hỏi aligned với LO
          \{lo\_code\} ở mức Bloom \{bloom\_level\}. Độ khó \{difficulty\}.'' 
\end{itemize}

\subsubsection{BGE-M3 Embedder}

Adapter \texttt{BGEEmbedder} chạy local model
\texttt{BAAI/bge-m3} qua thư viện \texttt{sentence-transformers},
sinh vector 1024 chiều. Hỗ trợ FP16 để giảm VRAM. Không phụ
thuộc API bên ngoài.

\subsection{Thiết kế Curriculum-Aware Layer (Điểm mới)}

Đây là điểm khác biệt chính của đồ án so với các hệ thống LMS
truyền thống --- tích hợp \textbf{Chuẩn đầu ra học phần (LO)}
vào toàn bộ pipeline sinh nội dung.

\subsubsection{DCMH Extractor (388 dòng)}

\texttt{DcmhExtractor} parse Đề cương môn học theo định dạng
CDIO của BKHCM bằng regex, không dùng LLM (structural-only,
đảm bảo deterministic):

\begin{itemize}
    \item Regex pattern cho LO: \texttt{r"L\textbackslash.O\textbackslash.
          \textbackslash d+\textbackslash.\textbackslash d+"}
    \item Regex cho chapter: \texttt{r"[Cc]h\textbackslash u\textbackslash
          o\textbackslash ng\textbackslash s+\textbackslash d+"}
    \item Bloom level detection: ánh xạ 40+ động từ
          tiếng Việt + tiếng Anh sang 6 mức Bloom.
    \item CDIO level: ánh xạ L.O.X.Y $\to$ level I (X=1,2),
          II (X=3,4), III (X=5,6).
\end{itemize}

\subsubsection{LO-Chunk Mapping}

\texttt{HeuristicLoMapper} tự động liên kết chunk với LO qua
heading alignment --- không cần LLM, nhanh và reproducible:

\begin{enumerate}
    \item Trích xuất chapter number từ heading (vd: ``Chương 4''
          $\to$ chapter 4).
    \item Map LO đến chapter: L.O.4.2 $\to$ chapter 4.
    \item Text matching: heading có chứa $\geq$ 2 từ khóa
          quan trọng từ statement\_vi của LO.
\end{enumerate}

Kết quả được persist vào bảng \texttt{chunk\_lo\_mappings} và
Neo4j relationship \texttt{SUPPORTS}.

\subsubsection{LO-Driven Quiz Generation}

\texttt{GenerateCurriculumQuizUseCase} gọi Gemini với prompt
specialized theo LO:

\begin{lstlisting}[language=Python, caption={Prompt template cho LO-aligned quiz}]
prompt = f"""You are an assessment designer.
Create {count} quiz items aligned with:
  LO Code: {lo.code}
  Statement: {lo.statement_vi}
  Bloom Level: {lo.bloom_level}
  Style: {style}  # quiz | midterm | final
Context from course materials:
{context_chunks}
"""
\end{lstlisting}

\subsection{Trạng thái triển khai MVP (snapshot 2026-05-10)}

Bảng~\ref{tab:gap-matrix} tổng hợp trạng thái hiện tại của hệ thống tính
đến ngày 10/05/2026, phân loại theo ba nhóm: \textbf{Đã triển khai},
\textbf{Đang hiện thực cho MVP GĐ1}, và \textbf{Chưa thuộc MVP GĐ1}.

\begin{table}[H]
    \centering
    \caption{Gap Matrix --- Trạng thái triển khai hệ thống (snapshot 2026-05-10).}
    \label{tab:gap-matrix}
    \begin{tabular}{p{5.5cm}p{7cm}c}
        \toprule
        \textbf{Thành phần} & \textbf{Mô tả} & \textbf{Trạng thái} \\
        \midrule
        \multicolumn{3}{c}{\textbf{Đã triển khai (Có sẵn trong codebase)}} \\
        \midrule
        gRPC Server + 12 RPCs
            & 8 RPCs cơ bản + 4 RPCs curriculum
            & \cellcolor{green!20}Hoàn thiện \\
        \midrule
        PostgreSQL (14 bảng)
            & 8 core + 6 curriculum tables
            & \cellcolor{green!20}Hoàn thiện \\
        \midrule
        DCMH Extractor (388 dòng)
            & Regex CDIO + Bloom level detection
            & \cellcolor{green!20}Hoàn thiện \\
        \midrule
        Neo4j Knowledge Graph
            & 4 curriculum methods
            & \cellcolor{green!20}Hoàn thiện \\
        \midrule
        HTTP API (FastAPI)
            & 11 endpoints + Swagger UI
            & \cellcolor{green!20}Hoàn thiện \\
        \midrule
        3 BullMQ Workers
            & document, enrichment, outbox
            & \cellcolor{green!20}Hoàn thiện \\
        \midrule
        Docker Compose Stack
            & 9 services infrastructure + app
            & \cellcolor{green!20}Hoàn thiện \\
        \midrule
        \multicolumn{3}{c}{\textbf{Đang hiện thực cho MVP GĐ1}} \\
        \midrule
        LMS Integration Schema
            & \texttt{lms\_user\_mappings} + \texttt{lms\_course\_ref}
            & \cellcolor{yellow!20}Đã thêm schema \\
        \midrule
        Endpoint \texttt{GET /api/documents}
            & Document listing cho LMS/web admin
            & \cellcolor{yellow!20}Đã thêm route \\
        \midrule
        HTTP API Service (compose)
            & \texttt{http-api} entry trong compose
            & \cellcolor{yellow!20}Đã thêm service \\
        \midrule
        Next.js Frontend
            & LTI/OIDC + cards + quiz + review
            & \cellcolor{yellow!20}Đã scaffold \\
        \midrule
        XBlock OIDC Fix
            & Sửa URL + OIDC flow
            & \cellcolor{yellow!20}Đã sửa \\
        \midrule
        \texttt{lms\_shared} network
            & Docker network cho 2 stack
            & \cellcolor{yellow!20}Đã config \\
        \midrule
        LTI Keypair + JWKS
            & RSA 2048 + endpoint
            & \cellcolor{yellow!20}Đã tạo \\
        \midrule
        \multicolumn{3}{c}{\textbf{Chưa thuộc MVP GĐ1}} \\
        \midrule
        LTI AGS Grade Passback
            & Tự động publish điểm về LMS
            & \cellcolor{red!20}GĐ2 \\
        \midrule
        Video Processing Pipeline
            & Whisper + scene detection
            & \cellcolor{red!20}GĐ2 \\
        \midrule
        RAG Chatbot
            & Retrieval-augmented Q\&A
            & \cellcolor{red!20}GĐ2 \\
        \midrule
        BKT + Mastery Tracking
            & Bayesian Knowledge Tracing
            & \cellcolor{red!20}GĐ2 \\
        \midrule
        Recommendation Engine
            & Personalized learning path
            & \cellcolor{red!20}GĐ2 \\
        \bottomrule
    \end{tabular}
\end{table}

\textbf{Ghi chú:} Backend \texttt{chunking\_v2} đã hoàn thiện phần lõi
với curriculum-aware layer từ giai đoạn đầu, vượt xa dự kiến ban đầu.
Điều này cho phép GĐ1 tập trung vào integration (LTI, Next.js frontend,
XBlock) thay vì phát triển backend từ đầu. Các tính năng GĐ2 (AGS
grade passback, video, chatbot, BKT, recommendation) sẽ được triển khai
trong giai đoạn sau và không nằm trong phạm vi MVP demo GĐ1.
